\documentclass[11pt,a4paper]{article}

\usepackage[utf8]{inputenc}
\usepackage[T1]{fontenc}
\usepackage[english]{babel}

\usepackage{helvet} 
\renewcommand{\familydefault}{\sfdefault}

\usepackage[left=2.5cm, right=2.5cm, top=1.5cm, bottom=1.5cm]{geometry}
\usepackage{parskip}

\usepackage{array}
\usepackage{xcolor}
\usepackage{enumitem}
\usepackage{graphicx}

\usepackage{fancyhdr}
\usepackage{lastpage}
\pagestyle{fancy}
\fancyhf{}
\renewcommand{\headrulewidth}{0.4pt}
\renewcommand{\footrulewidth}{0pt}

% Textos del encabezado y pie traducidos
\lhead{\textcolor{gray}{\small Document Title}} 
\rhead{\textcolor{gray}{\small \today}} 
\cfoot{\small Page \thepage\ of \pageref*{LastPage}} 

\usepackage[colorlinks=true, linkcolor=blue, urlcolor=blue]{hyperref}

\begin{document}

\section*{Introduction}
This is your improved template. Thanks to the \texttt{parskip} package, paragraphs no longer have indentation; instead, they are separated by a vertical space, giving it a much cleaner look.

\section*{Lists and Styles}
By using \texttt{enumitem} and \texttt{xcolor}, you can quickly create stylized lists:
\begin{itemize}[leftmargin=*]
    \item First item of the list.
    \item Second item with an \href{https://en.wikipedia.org}{interactive blue link}.
\end{itemize}

\end{document}